\documentclass[a4paper,11pt]{report}
\usepackage[showexo=true,showcorr=false]{../packages/coursclasse}
%Commenter ou enlever le commentaire sur la ligne suivante pour montrer le niveau
\toggletrue{montrerNiveaux}
%permet de gérer l'espacement entre les items des env enumerate et enumitem
\usepackage{enumitem}
\setlist[enumerate]{align=left,leftmargin=1cm,itemsep=10pt,parsep=0pt,topsep=0pt,rightmargin=0.5cm}
\setlist[itemize]{align=left,labelsep=1em,leftmargin=*,itemsep=0pt,parsep=0pt,topsep=0pt,rightmargin=0cm}
%permet de gerer l'espacement entre les colonnes de multicols
\setlength\columnsep{35pt}


\begin{document}
%%%%%%%%%%%%%%%%% À MODIFIER POUR CHAQUE SERIE %%%%%%%%%%%%%%%%%%%%%%%%%%%%%
\newcommand{\chapterName}{Nombres et Opérations}
\newcommand{\serieName}{Opérations avec des nombres relatifs}


%%%%%%%%%%%%%%%%%% PREMIERE PAGE NE PAS MODIFER %%%%%%%%%%%%%%%%%%%%%%%%
% le chapitre en cours, ne pas changer au cours d'une série
\chapter*{\chapterName}
\thispagestyle{empty}

%%%%% LISTE AIDE MEMOIRE %%%%%%
\begin{amL}{\serieName}{
\item Distance à zéro (page 17)
\item Nombres opposés (page 17)
\item Addition de nombres relatifs (page 18)
\item Soustraction de nombres relatifs (page 18)
\item Ecriture simplifiée d'une somme ou d'une différence de nombres relatifs (page 19)
\item Effectuer un calcul comportant des sommes et des différences de nombres relatifs (page 19)
}
\end{amL}
%%%%%%%%%%%%%%% DEBUT DE LA SERIE NE PAS MODIFIER %%%%%%%%%%%%%%%%%%%%%%%%%%%%%
\section*{\serieName}
\setcounter{page}{1}
\thispagestyle{firstPage}



%%%%%%%%%%% LES EXERCICES %%%%%%%%%%%%%%%%%%%%%%%%%%%%%%%%%%%


\begin{resolu}{Un jeu de billes}
{
\newcommand{\bverte}{\includegraphics[scale=.04]{media/no-37/smileyHappy.pdf}}
\newcommand{\brouge}{\includegraphics[scale=.04]{media/no-37/smileyAngry.pdf}}
Il y a deux armées: les billes positives \bverte et les billes négatives \brouge .
\begin{itemize}
\item[$\bullet$] Les billes d'une  même armée se rassemblent.
\item[$\bullet$] Deux billes d'armées différentes se neutralisent.
\end{itemize}
\begin{tasks}
\task Traduis les situations suivantes par un calcul puis trouve la réponse.

\vspace*{.5cm}

\bverte\bverte\bverte\bverte$\quad$ rencontrent $\quad$\brouge\brouge\brouge \hfill $(+4)+(-3) = {\color{blue}(+1)}$

\vspace*{.5cm}

\bverte\bverte$\quad$ rencontrent $\quad$\brouge\brouge\brouge\brouge \hfill $(+2)+(-4) = {\color{blue}(-2)}$

\vspace*{.5cm}

\brouge\brouge\brouge\brouge\brouge$\quad$ rencontrent $\quad$\bverte\bverte \hfill $({\color{blue}-5})+({\color{blue}+2}) = {\color{blue}(-3)}$

\vspace*{.5cm}

\brouge\brouge\brouge$\quad$ rencontrent $\quad$ \brouge\brouge\brouge \hfill $({\color{blue}-3})+({\color{blue}-3}) = {\color{blue}(-6)}$

\vspace*{.5cm}

\bverte\bverte\bverte$\quad$ rencontrent $\quad$\brouge\brouge\brouge \hfill $({\color{blue}+3})+({\color{blue}-3}) = {\color{blue}0}$

\vspace*{.5cm}

\brouge\brouge$\quad$ rencontrent $\quad$ \bverte\bverte\bverte\bverte\bverte \hfill $({\color{blue}-2})+({\color{blue}+5}) = {\color{blue}+3}$

\vspace*{.5cm}

7 positifs rencontrent 20 négatifs \hfill $({\color{blue}+7})+({\color{blue}-20}) = {\color{blue}-13}$

\vspace*{.5cm}

\task Calcule les sommes suivantes:
\begin{multicols}{2}
\begin{enumerate}
\item $(+3)+(+8)={\color{blue} +11 }$
\item $(-2)+(+3)={\color{blue}  +1}$
\item $(-2)+(-4)={\color{blue}  -6}$
\item $(+7)+(+10)={\color{blue} +17 }$
\item $(-1)+(-4)= {\color{blue}  -5}$
\item $(+7)+(-5)={\color{blue}  +2}$
\item $(+5)+(-1)={\color{blue}  +4}$
\item $(+15)+(+3)={\color{blue}  +18}$
\item $(-6)+(-5)={\color{blue}  -11}$
\item $(-3)+(+1)={\color{blue}  -2}$
\item $(+4)+(-6)={\color{blue}  -2}$
\item $(-2)+(+8)={\color{blue}  +6}$
\end{enumerate}
\end{multicols}
\item Complète le texte.
\begin{enumerate}
\item[$\bullet$] Pour additionner des nombres de même signe:
\begin{enumerate}
\item[$\rightarrow$] on {\color{blue}additionne} les distances à zéro des nombres~;
\item[$\rightarrow$] on met devant le résultat le signe {\color{blue}d'un des deux nombres}.
\end{enumerate}
\item[$\bullet$] Pour additionner des nombres de signes différents~:
\begin{enumerate}
\item[$\rightarrow$] on {\color{blue}soustrait} la plus petite distance à zéro à la plus grande;
\item[$\rightarrow$] on met devant le résultat le signe {\color{blue}du nombre qui a la plus grande distance à zéro}.
\end{enumerate}
\end{enumerate}
\end{tasks}
}{2}
\end{resolu}

\begin{exop}
{Additionne les nombres relatifs de mêmes signes.
\begin{tasks}(2)
\task $(-9)+(-16)=$
\task $(-28)+(-4)=$
\task $(+15)+(+12)=$
\task $(-45)+(-14)=$
\end{tasks}}{2}
\end{exop}

\begin{exop}
{Additionne les nombres relatifs de signes différents en répondant aux questions ci-dessous.
\begin{tasks}[after-item-skip=0.2em]
\task $(-15)+(+3)$

Dans l'addition ci-dessus,

Quel est le signe du nombre qui a la plus grande distance à zéro~? \hrulefill

Quel sera le signe du résultat~? \hrulefill

Quelle opération dois-je faire avec le 15 et le 3~? \hrulefill

Donc $(-15)+(+3)=$ \hrulefill 

\task $(+17)+(-20)$

Dans l'addition ci-dessus,

Quel est le signe du nombre qui a la plus grande distance à zéro~? \hrulefill

Quel sera le signe du résultat~? \hrulefill

Quelle opération dois-je faire avec le 17 et le 20~? \hrulefill

Donc $(+17)+(-20)=$ \hrulefill 

\task $(-11)+(+18)$

Dans l'addition ci-dessus,

Quel est le signe du nombre qui a la plus grande distance à zéro~? \hrulefill

Quel sera le signe du résultat~? \hrulefill

Quelle opération dois-je faire avec le 11 et le 18~? \hrulefill

Donc $(-11)+(+18)=$ \hrulefill 

\task $(-17)+(+9)$

Dans l'addition ci-dessus,

Quel est le signe du nombre qui a la plus grande distance à zéro~? \hrulefill

Quel sera le signe du résultat~? \hrulefill

Quelle opération dois-je faire avec le 17 et le 9~? \hrulefill

Donc $(-17)+(+9)=$ \hrulefill 

\end{tasks}

}{2}
\end{exop}

\begin{exop}{
Additionne les nombres relatifs de signes différents suivants.
\begin{tasks}(2)
\task $(+13)+(-11)=$
\task $(-10)+(+9)=$
\task $(-53)+(+62)=$
\task $(+15)+(-5)=$
\task $(+21)+(-15)=$
\end{tasks}}{2}
\end{exop}

\begin{exop}
{
Additionne les nombres relatifs suivants en répondant aux questions ci-dessous.
\begin{tasks}[after-item-skip=0.2em]
\task $(-15)+(-3)$

Dans l'addition ci-dessous,

Les deux nombres sont-ils de même signe~? \hrulefill

Si non, quel est le signe du nombre qui a la plus grande distance à 0~? \hrulefill

Quel sera le signe du résultat~? \hrulefill

Quelle opération dois-je faire avec le 15 et le 3~? \hrulefill

Donc $(-15)+(-3)=$ \hrulefill

\task $(-16)+(+4)$

Dans l'addition ci-dessous,

Les deux nombres sont-ils de même signe~? \hrulefill

Si non, quel est le signe du nombre qui a la plus grande distance à 0~? \hrulefill

Quel sera le signe du résultat~? \hrulefill

Quelle opération dois-je faire avec le 16 et le 4~? \hrulefill

Donc $(-16)+(+4)=$ \hrulefill

\task $(-4)+(+12)$

Dans l'addition ci-dessous,

Les deux nombres sont-ils de même signe~? \hrulefill

Si non, quel est le signe du nombre qui a la plus grande distance à 0~? \hrulefill

Quel sera le signe du résultat~? \hrulefill

Quelle opération dois-je faire avec le 4 et le 12~? \hrulefill

Donc $(-4)+(+12)=$ \hrulefill

\task $(-25)+(-18)$

Dans l'addition ci-dessous,

Les deux nombres sont-ils de même signe~? \hrulefill

Si non, quel est le signe du nombre qui a la plus grande distance à 0~? \hrulefill

Quel sera le signe du résultat~? \hrulefill

Quelle opération dois-je faire avec le 25 et le 18~? \hrulefill

Donc $(-25)+(-18)=$ \hrulefill

\task $(+8)+(-25)$

Dans l'addition ci-dessous,

Les deux nombres sont-ils de même signe~? \hrulefill

Si non, quel est le signe du nombre qui a la plus grande distance à 0~? \hrulefill

Quel sera le signe du résultat~? \hrulefill

Quelle opération dois-je faire avec le 8 et le 25~? \hrulefill

Donc $(+8)+(-25)=$ \hrulefill

\task $(-40)+(+52)$

Dans l'addition ci-dessous,

Les deux nombres sont-ils de même signe~? \hrulefill

Si non, quel est le signe du nombre qui a la plus grande distance à 0~? \hrulefill

Quel sera le signe du résultat~?

Quelle opération dois-je faire avec le 40 et le 52~? \hrulefill

Donc $(-40)+(+52)=$ \hrulefill
\end{tasks}
}{2}
\end{exop}

\begin{exop}{
Complète les égalités suivantes avec le nombre qui convient:
\begin{tasks}(2)
\task $(+7)+(\ligne{1.5})=(+11)$
\task $(+8)+(\ligne{1.5})=(+5)$
\task $(+13)+(\ligne{1.5})=0$
\task $(+5)+(\ligne{1.5})=(-3)$
\task $(-7)+(\ligne{1.5})=(+11)$
\task $(-13)+(\ligne{1.5})=0$
\task $(-4)+(\ligne{1.5})=(-9)$
\task $(+3)+(\ligne{1.5})=(-3)$
\end{tasks}
}{2}\end{exop}

\begin{resolu}
{Températures}{Résous les problèmes suivants.
\begin{tasks}
	\task La température relevée un matin d'hiver est de \tunit{-8}{\degreeCelsius}. Le lendemain, la température a augmenté de \tunit{5}{\degreeCelsius}. Quelle est alors cette température~?

{\bf Réponse :} On traduit l'énoncé par le calcul suivant : $(-8)+(+5)=(-3).$ 

La température du lendemain est de \tunit{-3}{\degreeCelsius}.

\task Ma cave fait \tunit{22}{\degreeCelsius} de plus que mon congélateur qui affiche \tunit{-20}{\degreeCelsius}. Elle a aussi \tunit{17}{\degreeCelsius} de moins que mon appartement. Quelle est la température de mon appartement~?

{\bf Réponse :} On traduit l'énoncé par les calculs suivants : $$(-20)+(+22)=(+2) \qquad (+2)+(+17)= (+19).$$

La température de mon appartement est de \tunit{19}{\degreeCelsius}.
\end{tasks}

}{2}
\end{resolu}

\begin{exo}
{Résous les problèmes suivants.
\begin{tasks}
\task La température relevée un matin d'hiver est de \tunit{-5}{\degreeCelsius}. Le lendemain, la température a augmenté de \tunit{3}{\degreeCelsius}. Quelle est alors cette température~?
\task Ma cave fait \tunit{20}{\degreeCelsius} de plus que mon congélateur qui affiche \tunit{-18}{\degreeCelsius}. Elle a aussi \tunit{15}{\degreeCelsius} de moins que mon appartement. Quelle est la température de mon appartement~?
\end{tasks}
}{2}
\end{exo}

\begin{resolu}{Regrouper}
{Calcule les expressions suivantes après avoir regroupé les positifs et les négatifs.
\begin{tasks}
\task $A=(-2)+(+3)+(-4)+(+5)=$

$A=(-2)+(-4)+(+3)+(+5)=(-6)+(+8)=(+2)$
\task $B=(-7)+(+9)+(-13)+(+11)=$

$B=(-7)+(-13)+(+9)+(+11)=(-20)+(+20)=0$
\task $C=(-15)+(-6)+(+9)+(-11)+(+17)=$

$C=(-15)+(-6)+(-11)+(+9)+(+17)=(-32)+(+26)=(-6)$
\end{tasks}
}{2}
\end{resolu}


\begin{exo}
{Calcule les expressions suivantes après avoir regroupé les positifs et les négatifs.
\begin{tasks}
\task $A=(-5)+(+8)+(-11)+(+13)$
\task $B=(-16)+(-4)+(+7)+(-13)+(+15)$
\task $C=(-1)+(+2)+(-3)+(+4)+(-5)$
\task $D=(-4,5)+(+11,\!7)+(+9,\!3)+(-5,\!8)$
\end{tasks}
}{2}
\end{exo}

\begin{exo}
{Peux-tu trouver $\ldots$
\begin{tasks}
\task $\ldots$ deux nombres négatifs dont la somme vaut $-13$~?
\task $\ldots$ deux nombres égaux dont la somme vaut $-16$~?
\task $\ldots$ deux nombres de signes contraires dont la somme vaut $-13$~?
\task $\ldots$ trois nombres relatifs dont la somme vaut $-13$~?
\task $\ldots$ deux nombres positifs dont la somme vaut $-13$~?
\end{tasks}
}{2}
\end{exo}

\exof{NO120}{42}{2}
\exof{NO121}{42}{3}

\begin{resolu}{Soustraction avec les relatifs}{
Observe les exemples suivants:

\bigskip

$\begin{array}{|ccccc|}\hline
(+8)&{\color{magenta}-}&({\color{red}+}5)&= & \\
& {\color{magenta}\downarrow}&{\color{red}\downarrow} & & \\
(+8)&{\color{magenta}+}&({\color{red}-}5)&=&(+3)\\\hline
\end{array}$ 
\hfill 
$\begin{array}{|ccccc|}\hline
(-8)&{\color{magenta}-}&({\color{red}+}5)&= & \\
& {\color{magenta}\downarrow}&{\color{red}\downarrow} & & \\
(-8)&{\color{magenta}+}&({\color{red}-}5)&=&(-13)\\\hline
\end{array}$ 

\bigskip

$\begin{array}{|ccccc|}\hline
(+8)&{\color{magenta}-}&({\color{red}-}5)&= & \\
& {\color{magenta}\downarrow}&{\color{red}\downarrow} & & \\
(+8)&{\color{magenta}+}&({\color{red}+}5)&=&(+13)\\\hline
\end{array}$ 
\hfill 
$\begin{array}{|ccccc|}\hline
(-8)&{\color{magenta}-}&({\color{red}-}5)&= & \\
& {\color{magenta}\downarrow}&{\color{red}\downarrow} & & \\
(-8)&{\color{magenta}+}&({\color{red}+}5)&=&(-3)\\\hline
\end{array}$ 

\bigskip

Complète le texte avec les mots proposés:

\begin{center}
\fbox{addition} $\quad$ \fbox{soustraction} $\quad$ \fbox{deuxième} $\quad$ \fbox{positif} $\quad$ \fbox{négatif} 
\end{center}

Pour soustraire deux nombres relatifs, il faut:\begin{itemize}
\item[{\color{magenta}$\rightarrow$}] {\color{magenta} Transformer la {\color{blue} soustraction} en {\color{blue} addition}.}
\item[{\color{red}$\rightarrow$}] {\color{red} Prendre l'opposé du {\color{blue} deuxième} terme.

C'est-à-dire qu'un terme {\color{blue} négatif} devient {\color{blue} positif} et inversement.}
\end{itemize}
}{2}
\end{resolu}


\begin{exop}
{Effectue les soustractions suivantes.
\begin{tasks}(2)
\task $\begin{array}{|ccccc|}\hline
(+24)&{\color{magenta}-}&({\color{red}-}17)&= & \\
& {\color{magenta}\downarrow}&{\color{red}\downarrow} & & \\
(\ligne{1})&{\color{magenta}\ligne{.5}}&(\ligne{1})&=&(\ligne{1})\\\hline
\end{array}$ 
\task $\begin{array}{|ccccc|}\hline
(+8)&{\color{magenta}-}&({\color{red}+}5)&= & \\
& {\color{magenta}\downarrow}&{\color{red}\downarrow} & & \\
(\ligne{1})&{\color{magenta}\ligne{.5}}&(\ligne{1})&=&(\ligne{1})\\\hline
\end{array}$ 
\task $\begin{array}{|ccccc|}\hline
(-18)&{\color{magenta}-}&({\color{red}+}29)&= & \\
& {\color{magenta}\downarrow}&{\color{red}\downarrow} & & \\
(\ligne{1})&{\color{magenta}\ligne{.5}}&(\ligne{1})&=&(\ligne{1})\\\hline
\end{array}$ 
\task $\begin{array}{|ccccc|}\hline
(-35)&{\color{magenta}-}&({\color{red}+}21)&= & \\
& {\color{magenta}\downarrow}&{\color{red}\downarrow} & & \\
(\ligne{1})&{\color{magenta}\ligne{.5}}&(\ligne{1})&=&(\ligne{1})\\\hline
\end{array}$ 
\end{tasks}
}{2}
\end{exop}

\begin{exop}
{Effectue les soustractions suivantes.
\begin{tasks}(2)
\task $\begin{array}{|ccccc|}\hline
(+5)&{\color{magenta}-}&({\color{red}-}1)&= & \\
(\ligne{1})&{\color{magenta}+}&(\ligne{1})&=&(\ligne{1})\\\hline
\end{array}$ 
\task $\begin{array}{|ccccc|}\hline
(-5)&{\color{magenta}-}&({\color{red}+}3)&= & \\
(\ligne{1})&{\color{magenta}+}&(\ligne{1})&=&(\ligne{1})\\\hline
\end{array}$ 
\task $\begin{array}{|ccccc|}\hline
(+4)&{\color{magenta}-}&({\color{red}-}8)&= & \\
(\ligne{1})&{\color{magenta}+}&(\ligne{1})&=&(\ligne{1})\\\hline
\end{array}$ 
\task $\begin{array}{|ccccc|}\hline
(+4)&{\color{magenta}-}&({\color{red}+}1)&= & \\
(\ligne{1})&{\color{magenta}+}&(\ligne{1})&=&(\ligne{1})\\\hline
\end{array}$ 
\end{tasks}
}{2}
\end{exop}

\begin{exop}
{Effectue les soustractions suivantes.
\begin{tasks}(2)
\task $(-5)-(+8)=$

\bigskip\ligne{4} $=$ \ligne{1}

\task $(+2)-(+2)=$

\bigskip\ligne{4} $=$ \ligne{1}

\task $(+10)-(+3)=$

\bigskip\ligne{4} $=$ \ligne{1}

\task $(-10)-(+6)=$

\bigskip\ligne{4} $=$ \ligne{1}

\task $(+8)-(-8)=$

\bigskip\ligne{4} $=$ \ligne{1}

\task $(-25)-(+35)=$

\bigskip\ligne{4} $=$ \ligne{1}

\task $(-7)-(+17)=$

\bigskip\ligne{4} $=$ \ligne{1}

\task $(-9)-(-8)=$

\bigskip\ligne{4} $=$ \ligne{1}

\end{tasks}
}{2}
\end{exop}

\begin{exo}
{Effectue les soustractions suivantes.
\begin{tasks}(2)
\task $(-9)-(+6)$
\task $(+7)-(-8)$
\task $(-12)-(+5)$
\task $(+13)-(+4)$
\task $(-25)-(-30)$
\end{tasks}}{2}
\end{exo}

\begin{exo}
{Effectue les soustractions suivantes.
\begin{tasks}(2)
\task $(+15,\!1)-(+18,\!1)$
\task $(-14,\!05)-(-19,\!9)$
\task $(-204)-(+170)$
\task $(-41,\!48)-(+42,\!28)$
\task $(+14)-(-12,\!14)$
\end{tasks}}{3}
\end{exo}

\begin{exop}
{Parmi les écritures suivantes encadre celles qui sont exactes:
\begin{tasks}(2)
\task $(+1673)-(+1622)=(+51)$
\task $(+121)-(+180)=(+59)$
\task $(-30)-(-69)=(+39)$
\task $(-69)-(-30)=(+39)$
\task $(-63)-(-14)=(+77)$
\task $(+14)-(-63)=(+77)$
\task $(-44)-(-101)=(+57)$
\task $(-10)-(+54)=(+64)$
\end{tasks}
}{3}
\end{exop}


\begin{exo}
{Calcule les chaines d'opérations suivantes:
\begin{tasks}(2)
\task $A=(+20)+(+11)+(-18)=$
\task $B=(-5)+(+6)-(-7)=$
\task $C=(+8)-(+12)+(-21)=$
\task $D=(+9)+(-12)-(-6)=$
\end{tasks}
}{2}
\end{exo}

\begin{exo}
{Calcule les chaines d'opérations suivantes:
\begin{tasks}
\task $A=(+6)-(+5)+(-9)+(+8)+(-10)-(-5)=$
\task $B=(+12)+(-15)-(-7)+(-27)-(+4)=$
\task $C=(-7)+(+5)-(+3)+(-8)-(+7)+(+5)=$
\task $D=(+2)-(-1)+(+9)-(+6)+(-3)=$
\end{tasks}
}{2}
\end{exo}

\begin{exo}
{Calcule les chaines d'opérations suivantes:
\begin{tasks}
\task $A=(-2,\!5)+(+6)-(-7)+(+2,\!5)+(+11)+(-18)=$
\task $B=(+8,\!3)-(-12)+(-21)+(+9,\!2)+(+12)-(-6)=$
\task $C=(-6,\!7)-(+5)+(-9,\!3)+(+3,\!8)+(-10,\!7)+(+5)=$
\task $D=(+0,\!12)+(-1,\!5)+(+0,\!7)+(-2,\!74)+(+1,\!24)=$
\end{tasks}
}{3}
\end{exo}

\exof{NO127}{44}{2}
\exof{NO128}{44}{2}
\exof{NO129}{44}{3}

\begin{resolu}{Simplification d'écriture}{Observe les exemples suivants:

\bigskip

$\begin{array}{|ccccc|}\hline
(+5)&{\color{olive}\ovalbox{\color{black}$+(+$}}3)&=&(+8)& \mbox{devient} \\
{\color{violet}\downarrow}&{\color{olive}\downarrow}& & {\color{violet}\downarrow} & \\
5 & +3 & = & 8 & \\\hline
\end{array}$
\hfill
$\begin{array}{|ccccc|}\hline
(-5)&{\color{olive}\ovalbox{\color{black}$+(-$}}3)&=&(-8)& \mbox{devient} \\
{\color{violet}\downarrow}&{\color{olive}\downarrow}& & {\color{violet}\downarrow} & \\
-5 & -3 & = & -8 & \\\hline
\end{array}$

$\begin{array}{|ccccc|}\hline
(+5)&{\color{olive}\ovalbox{\color{black}$+(-$}}3)&=&(+2)& \mbox{devient} \\
{\color{violet}\downarrow}&{\color{olive}\downarrow}& & {\color{violet}\downarrow} & \\
5 & -3 & = & 2 & \\\hline
\end{array}$
\hfill
$\begin{array}{|ccccc|}\hline
(-5)&{\color{olive}\ovalbox{\color{black}$-(+$}}3)&=&(-8)& \mbox{devient} \\
{\color{violet}\downarrow}&{\color{olive}\downarrow}& & {\color{violet}\downarrow} & \\
-5 & -3 & = & -8 & \\\hline
\end{array}$

$\begin{array}{|ccccc|}\hline
(-5)&{\color{olive}\ovalbox{\color{black}$+(+$}}3)&=&(-2)& \mbox{devient} \\
{\color{violet}\downarrow}&{\color{olive}\downarrow}& & {\color{violet}\downarrow} & \\
-5 & +3 & = & -2 & \\\hline
\end{array}$
\hfill
$\begin{array}{|ccccc|}\hline
(-5)&{\color{olive}\ovalbox{\color{black}$-(-$}}3)&=&(-2)& \mbox{devient} \\
{\color{violet}\downarrow}&{\color{olive}\downarrow}& & {\color{violet}\downarrow} & \\
-5 & +3 & = & -2 & \\\hline
\end{array}$

\bigskip

Complète le texte avec les mots proposés:

\begin{center}
\fbox{signe} $\quad$ \fbox{plus} $\quad$ \fbox{positif} $\quad$ \fbox{négatif} $\quad$ \fbox{moins} $\quad$ \fbox{parenthèses} 
\end{center}

Pour simplifier une écriture de nombres relatifs, il faut:

\begin{tasks}
\task[{\color{black}$\rightarrow$}] Enlever toutes les {\color{blue}parathèses}.
\task[{\color{violet}$\rightarrow$}] {\color{violet} Si le premier nombre est {\color{blue} positif}, on enlève le {\color{blue} signe plus}}.
\task[{\color{violet}$\rightarrow$}] {\color{violet} Si le premier nombre est {\color{blue} négatif}, on laisse le {\color{blue} signe moins}}.
\task[{\color{olive}$\rightarrow$}] {\color{olive}Deux signes consécutifs sont remplacés par un seul en suivant la règle suivante:
\begin{align*}
++\quad\quad   \text{ et }\quad \quad -- \quad\quad \implies \quad \quad {\color{blue} +}\\
+-\quad\quad   \text{ et }\quad \quad -+ \quad\quad \implies \quad \quad {\color{blue} -}
\end{align*}}
\end{tasks}

}{3}
\end{resolu}

\begin{exop}
{Simplifie les écritures suivantes, puis effectue le calcul.

\bigskip

\begin{tasks}(2)
\task $\begin{array}{cccc}
(+17)&{\color{olive}\ovalbox{\color{black}$+(-$}}19)& \mbox{devient} \\
{\color{violet}\downarrow}&{\color{olive}\downarrow}& & \\
\end{array}$

\vspace*{.7cm}\ligne{3.2} $=$\ligne{1}

\task $\begin{array}{cccc}
(+10)&{\color{olive}\ovalbox{\color{black}$-(-$}}10)& \mbox{devient} \\
{\color{violet}\downarrow}&{\color{olive}\downarrow}& & \\
\end{array}$

\vspace*{.7cm}\ligne{3.2} $=$\ligne{1}

\task $\begin{array}{cccc}
(-19)&{\color{olive}\ovalbox{\color{black}$-(+$}}3)& \mbox{devient} \\
{\color{violet}\downarrow}&{\color{olive}\downarrow}& & \\
\end{array}$

\vspace*{.7cm}\ligne{3.2} $=$\ligne{1}

\task $\begin{array}{cccc}
(+36)&{\color{olive}\ovalbox{\color{black}$+(+$}}14)& \mbox{devient} \\
{\color{violet}\downarrow}&{\color{olive}\downarrow}& & \\
\end{array}$

\vspace*{.7cm}\ligne{3.2} $=$\ligne{1}

\task $\begin{array}{cccc}
(-18)&{\color{olive}\ovalbox{\color{black}$-(+$}}1)& \mbox{devient} \\
{\color{violet}\downarrow}&{\color{olive}\downarrow}& & \\
\end{array}$

\vspace*{.7cm}\ligne{3.2} $=$\ligne{1}

\task $\begin{array}{cccc}
(-30)&{\color{olive}\ovalbox{\color{black}$-(+$}}21)& \mbox{devient} \\
{\color{violet}\downarrow}&{\color{olive}\downarrow}& & \\
\end{array}$


\vspace*{.7cm}\ligne{3.2} $=$\ligne{1}

\end{tasks}

}{2}
\end{exop}


\begin{exop}
{Simplifie les écritures suivantes. Puis effectue le calcul.

\begin{tasks}(2)
\task $(+13)+(-15)$ devient

\vspace*{.7cm}\ligne{3.2} $=$\ligne{1}

\task $(+9)-(-3)$ devient

\vspace*{.7cm}\ligne{3.2} $=$\ligne{1}

\task $(-17)-(+23)$ devient

\vspace*{.7cm}\ligne{3.2} $=$\ligne{1}

\task $(+35)+(+24)$ devient

\vspace*{.7cm}\ligne{3.2} $=$\ligne{1}

\task $(-14)-(+2)$ devient

\vspace*{.7cm}\ligne{3.2} $=$\ligne{1}

\task $(+40)-(+21)$ devient

\vspace*{.7cm}\ligne{3.2} $=$\ligne{1}


\end{tasks}
}{2}
\end{exop}

\begin{exo}
{Simplifie l'écriture, puis calcule:

\begin{tasks}
\task $(+3)+(-6)-(+4)+(-7)-(-6)+(-3)$
\task $(+2)+(-5)-(-3)+(-4)+(+6)+(-3)$
\task $(-12)+(+27)+(-5)-(-4)+(+12)-(-17)$
\task $(-6)+(-12)-(+3)+(-4)-(-5)-(+3)$
\task $(+14)+(-15)-(+14)+(-6)-(-3)+(+15)$
\task $(+3)+(-12)-(+4)+(-6)-(-7)-(+4)$
\end{tasks}
}{2}
\end{exo}

\begin{exop}
{Parmi les écitures suivantes, entoure celle(s) qui représente(nt) le même nombre que $A$~?

$$A=(-3)+(-5)-(-5)+(+2)$$

\begin{tasks}(1)
	\task[] $B=-3-5-(4-5)+2$

	\task[] $C=-(+3)+(-5)-(+4)+(-5)+(+2)$

	\task[] $D=-3-(+5)-4-(+5)-(-2)$

	\task[] $E=(-3)-(-5)+(-4)-5+(+2)$
\end{tasks}
}{3}
\end{exop}


\exof{NO132}{46}{2}
\exof{NO133}{46}{2}
\exof{NO130}{45}{2}
\exof{NO135}{47}{2}
\exof{NO136}{47}{2}
\exof{NO137}{48}{2}
\exol{NO139}{41}{2}
\exol{NO140}{42}{2}

	\begin{qmoodle}{Addition et soustraction de relatifs -- deux termes}{3}{
	\begin{center}	
		Q-1

\includegraphics[scale=1]{media/qr/nerop1}

\tiny{{https://edu.ge.ch/qr/nerop1}}
\end{center}
	\begin{center}	
		Q-2

\includegraphics[scale=1]{media/qr/nerop2}

\tiny{{https://edu.ge.ch/qr/nerop2}}
\end{center}
	\begin{center}	
		Q-3

\includegraphics[scale=1]{media/qr/nerop3}

\tiny{{https://edu.ge.ch/qr/nerop3}}
\end{center}
}
\end{qmoodle}

\begin{qmun}{Addition et soustraction de relatifs}{
		\begin{center}
\includegraphics[scale=1]{media/qr/nerop4}

\tiny{{https://edu.ge.ch/qr/nerop4}}
		\end{center}
	}
\end{qmun}

\end{document}
